% FOCAL-UMS Method, revised against the executable implementation.
% Intended for Springer LNCS via % FOCAL-UMS Method, revised against the executable implementation.
% Intended for Springer LNCS via % FOCAL-UMS Method, revised against the executable implementation.
% Intended for Springer LNCS via % FOCAL-UMS Method, revised against the executable implementation.
% Intended for Springer LNCS via \input{method_revised}.
% Required packages in the main file: amsmath, amssymb, graphicx.

\section{Method}
\label{sec:method}

\subsection{Problem Setup and Framework Overview}
\label{sec:setup}

We consider exemplar-free multi-domain continual learning over a sequence of
$T$ datasets
$\mathcal{D}_{1:T}=\{\mathcal{D}_1,\ldots,\mathcal{D}_T\}$, where
$\mathcal{D}_t=\{(x_i^t,y_i^t)\}_{i=1}^{N_t}$ has a task-specific label
space $\mathcal{Y}_t$.  At step $t$, only $\mathcal{D}_t$ is available for
optimization; images from $\mathcal{D}_{<t}$ are neither stored nor replayed.
The evaluated sequence contains ten heterogeneous visual datasets and therefore
induces both pronounced domain shifts and different prediction heads.

The implementation instantiates two separately pretrained and frozen ViT-B/16
networks~\cite{dosovitskiy2021vit}.  The local backbone
$f_{\mathrm{loc}}$ is used for classification and receives additive modulation
in selected attention projections.  The routing backbone $f_{\mathrm{route}}$
remains unmodulated and provides a class-token feature for task selection.  For
each task $t$, a linear bottleneck $b_t:\mathbb{R}^{768}\rightarrow
\mathbb{R}^{500}$ and a classifier
$c_t:\mathbb{R}^{500}\rightarrow\mathbb{R}^{|\mathcal{Y}_t|}$ are learned and
then frozen.

FOCAL-UMS follows an encode--decode process.  Task-specific Fourier
coefficients first generate additive attention modulations.  After training,
each modulation is encoded as a fixed-length task coordinate in an incrementally
constructed \emph{Unified Modulation Subspace} (UMS); at inference, that
coordinate is decoded into a modulation matrix.  Subspace-aware training and
post-task refinement connect these two stages by exposing the task parameters
to the candidate representation during learning and by fitting the shared
rotation and support decoder after encoding.  Sparse Fourier parameterization
and Mahalanobis routing are adopted mechanisms.  The methodological focus of
FOCAL-UMS is the integration of UMS encoding and decoding into sequential
learning; the refinement objective is one component of this process rather than
a stand-alone collection of loss functions.

\begin{figure}[t]
  \centering
  \includegraphics[width=\linewidth]{\detokenize{FOCAL-UMS pipeline with post-task refinement.png}}
  \caption{FOCAL-UMS.  Task-specific Fourier coefficients are optimized through
  a candidate UMS map, encoded after training, and refined before decoding.  At
  inference, one task is selected for the complete grouped query set and its
  UMS code is decoded into attention modulations.  The anchor regularizes old
  rotation directions; it does not imply invariant old-task modulations.}
  \label{fig:focal_ums}
\end{figure}

\subsection{Task-Specific Frequency Modulation}
\label{sec:frequency}

We adapt sparse Fourier-domain fine-tuning~\cite{gao2024fourierft} to allocate
separate trainable coordinates to successive tasks.  Let $d=768$ and let $\pi$
be a deterministic, seeded permutation of the $d^2$ entries of a square
spectrum.  The first $Tn_f$ entries are partitioned into pairwise-disjoint
supports $\Omega_1,\ldots,\Omega_T$, each of cardinality $n_f$.  For task $t$,
the vector $\boldsymbol{\theta}_t\in\mathbb{R}^{n_f}$ is scattered onto
$\Omega_t$ to form
\begin{equation}
 S_t[u_m,v_m]=\theta_{t,m}\quad ((u_m,v_m)\in\Omega_t),
 \qquad S_t[u,v]=0\quad ((u,v)\notin\Omega_t).
 \label{eq:sparse_spectrum}
\end{equation}
The corresponding additive weight modulation is
\begin{equation}
 \Delta W_t=\alpha\,\operatorname{Re}\!\left(\mathcal{F}^{-1}_2(S_t)\right),
 \label{eq:ifft_modulation}
\end{equation}
where $\mathcal{F}^{-1}_2$ follows the default normalization of
\texttt{torch.fft.ifft2}.  In the primary configuration, $n_f=5{,}898$ and
$\alpha=300$.  The model modulates the query, key, value, and output projections
in blocks $6$--$11$, yielding 24 independent FocalExpert/UMS modules.  Their
pretrained projection weights remain frozen.

Disjoint supports prevent two tasks from directly updating the same allocated
raw spectral coordinates.  This statement is deliberately local to the raw
parameterization: taking the real part in Eq.~\eqref{eq:ifft_modulation}, and
later applying UMS reconstruction and support gating, does not establish exact
orthogonality of the resulting spatial updates or guarantee the absence of
forgetting.

\subsection{UMS-in-the-Loop Training}
\label{sec:inloop}

Each modulated projection owns a UMS with column bases
$P,Q\in\mathbb{R}^{d\times R_{\max}}$, a current active rank $r_t$, and a
shared rotation $O\in\mathbb{R}^{d\times d}$.  UMS-in-the-loop is designed to
reduce, rather than eliminate, the mismatch between unconstrained task training
and subsequent subspace encoding.  At task start and at periodic refreshes, the
current spectrum is transformed to the canonical frame,
\begin{equation}
 T_t=O^{\mathsf T}S_tO.
 \label{eq:canonical}
\end{equation}
The implementation removes the component explained by the active left basis,
computes a truncated SVD of the residual, residualizes the resulting right
directions against $Q$, and orthonormalizes both sides by QR.  Appending these
directions to the active columns gives candidate bases
$P_{\mathrm{c}},Q_{\mathrm{c}}\in\mathbb{R}^{d\times k}$, capped by
$R_{\max}$.

Let $A=OP_{\mathrm{c}}$ and $B=OQ_{\mathrm{c}}$.  Restricting the induced
bilinear representation to the current support gives
\begin{equation}
 G_t[m,j]=A[u_m,j]B[v_m,j],
 \qquad G_t\in\mathbb{R}^{n_f\times k}.
 \label{eq:G}
\end{equation}
The coefficients used in the forward pass are
\begin{equation}
 \widetilde{\boldsymbol{\theta}}_t
 =G_t\bigl(G_t^{\mathsf T}\boldsymbol{\theta}_t\bigr),
 \label{eq:inloop_map}
\end{equation}
and are scattered onto $\Omega_t$ before Eq.~\eqref{eq:ifft_modulation}.  The
matrix $G_t$ is cached without gradients and is refreshed at task start and
periodically during training.  Because its support-restricted columns are not
explicitly orthonormal and the code does not apply
$(G_t^{\mathsf T}G_t)^{-1}$, Eq.~\eqref{eq:inloop_map} is a
candidate-subspace-induced linear map, not an orthogonal projector.  Cross
entropy updates the current raw coefficients, bottleneck, and classifier.

\subsection{Post-Task UMS Encoding}
\label{sec:encoding}

The purpose of post-task encoding is to express each learned modulation through
shared bases and a task-specific coordinate that can later be decoded.  At task
completion, the current implementation rebuilds $S_t$ from the
\emph{raw} slice $\boldsymbol{\theta}_t$ and encodes this matrix into UMS.  It
does not fold the final in-loop value
$\widetilde{\boldsymbol{\theta}}_t$ from Eq.~\eqref{eq:inloop_map}; thus the
in-loop mechanism biases learning toward the candidate representation but does
not make training and end-task encoding identical.

Let $r_{t-1}$ be the active rank before encoding task $t$, with $r_0=0$.
For the previously active bases $P_{t-1}=P_{:,1:r_{t-1}}$ and
$Q_{t-1}=Q_{:,1:r_{t-1}}$, the left residual is
\begin{equation}
 U_t=T_t-P_{t-1}(P_{t-1}^{\mathsf T}T_t).
 \label{eq:left_residual}
\end{equation}
A truncated SVD
$U_t\approx U_s\operatorname{diag}(\boldsymbol{\sigma})V_s^{\mathsf T}$
provides paired residual directions.  The right directions are additionally
residualized against $Q_{t-1}$; QR-orthonormalized directions are appended to $P$
and $Q$ when
\begin{equation}
 \mathcal{E}_t=\|U_s\operatorname{diag}(\boldsymbol{\sigma})\|_F^2>\tau
 \label{eq:residual_energy}
\end{equation}
and capacity remains.  In the evaluated configuration,
$R_{\max}=768$, the truncated-SVD budget is 75, and $\tau=10^{-3}$.  With the
updated active rank $r_t$ and bases $P_r=P_{:,1:r_t}$ and
$Q_r=Q_{:,1:r_t}$, task $t$ is encoded as the zero-padded diagonal core
\begin{equation}
 z_t=\operatorname{pad}_{R_{\max}}\!\left[
 \operatorname{diag}\!\left(P_r^{\mathsf T}T_tQ_r\right)\right].
 \label{eq:task_code}
\end{equation}
This is a rank-budgeted diagonal approximation, not an exact factorization of
the complete core.  UMS codes are used for inference, although the present
training/checkpoint implementation also retains the preallocated raw
coefficient tensor.

\subsection{UMS Refinement and Decoding}
\label{sec:multiobjective}

The purpose of refinement is to reduce the discrepancy between the sparse task
spectrum and the spectrum decoded from the diagonal UMS coordinate.  A single
reconstruction term is insufficient for this purpose: it does not identify the
sparse support, expose off-diagonal energy omitted by $z_t$, or control changes
to previously used rotation directions.  We therefore optimize the rotation
and support decoder with a structured objective derived from these encoding
errors.

The shared orientation is parameterized by the Cayley transform
\begin{equation}
 A_{\mathrm{rot}}=W_{\mathrm{rot}}-W_{\mathrm{rot}}^{\mathsf T},\qquad
 O=\left(I-\tfrac{1}{2}A_{\mathrm{rot}}\right)^{-1}
   \left(I+\tfrac{1}{2}A_{\mathrm{rot}}\right),
 \label{eq:cayley}
\end{equation}
which is orthogonal in exact arithmetic whenever the solve is well-defined.
For the current code $z_t$, define the ungated spectral reconstruction and the
support logits as
\begin{align}
 \bar S_t &= OP\operatorname{diag}(z_t)Q^{\mathsf T}O^{\mathsf T},
 \label{eq:ungated_reconstruction}\\
 H_t &= W_L\operatorname{diag}(z_t)W_R^{\mathsf T}+\Gamma.
 \label{eq:support_logits}
\end{align}
Here $W_L,W_R\in\mathbb{R}^{d\times R_{\max}}$ and
$\Gamma\in\mathbb{R}^{d\times d}$ are shared support-decoder parameters.  Let
$\mathcal{I}_t=\{(u,v):S_t[u,v]\neq0\}$,
$Y_t[u,v]=\mathbf{1}[(u,v)\in\mathcal{I}_t]$, and
$C_t=P_r^{\mathsf T}(O^{\mathsf T}S_tO)Q_r$.  For nonempty $\mathcal{I}_t$,
let $m_t$ be the mean of $\bar S_t$ on $\mathcal{I}_t$ and let
$s_t=\sqrt{\operatorname{Var}_{\mathcal{I}_t}(\bar S_t)+10^{-8}}$,
where the variance uses the population denominator.  The code $z_t$ is held
fixed during refinement.  For $t>1$, let $O^{-},P^{-},Q^{-}$ denote the
pre-refinement rotation and previously active basis columns, with
$r_{t-1}$ columns.  The effective refinement objective is
\begin{align}
 \mathcal{L}_{\mathrm{UMS}}
 &=\lambda_{\mathrm{val}}
   \frac{\sum_{(u,v)\in\mathcal{I}_t}(S_t[u,v]-\bar S_t[u,v])^2}
        {|\mathcal{I}_t|}
 +\lambda_{\mathrm{mask}}\operatorname{BCEWithLogits}(H_t,Y_t;w_+)
 \notag\\
 &\quad+\lambda_{\mathrm{mom}}\bigl(m_t^2+(s_t-1)^2\bigr)
 +\lambda_{\mathrm{off}}
   \frac{\|C_t-\operatorname{diag}(\operatorname{diag}C_t)\|_F^2}
        {\max\{r_t(r_t-1),1\}}
 \notag\\
 &\quad+\lambda_{\mathrm{drift}}
   \frac{\|(O-O^{-})P^{-}\|_F^2+\|(O-O^{-})Q^{-}\|_F^2}
        {2r_{t-1}},
 \label{eq:loss_total}
\end{align}
where $w_+=(d^2-n_f)/n_f$ compensates for support imbalance.  The first two
terms penalize value and support-mask errors, respectively; the moment term
regularizes the scale of decoded occupied entries; the off-diagonal term
penalizes canonical-core energy omitted by the diagonal code; and the drift
term penalizes rotation changes along previously active basis directions.  The
value and moment terms are zero when $\mathcal{I}_t$ is empty; the drift term
is zero when there are no previously active columns.  The implemented weights are
$(\lambda_{\mathrm{val}},\lambda_{\mathrm{mask}},\lambda_{\mathrm{mom}},
\lambda_{\mathrm{off}},\lambda_{\mathrm{drift}})
=(1.0,0.1,0.01,0.1,2.0)$.  Refinement uses 15 Adam steps with learning rates
$10^{-4}$ for $W_{\mathrm{rot}}$ and $10^{-3}$ for the support decoder, after
which the current task is re-encoded without further basis expansion.

The implementation also forms a nominal drift penalty for the previously
active columns of $W_L$ and $W_R$.  Its reference tensors are detached views,
however, rather than independent copies: they share storage with the current
parameters and change when Adam updates them.  Consequently, their difference
from the current columns and the resulting penalty and gradient are zero.
Equation~\eqref{eq:loss_total} therefore states the effective drift term.  The
anchor does not constrain $\Gamma$, re-encode all previous tasks, or establish
invariant old-task readouts.  Numerical orthogonality of the Cayley rotation
can be assessed using the saved diagnostics.

At readout, the support decoder produces a hard gate and the task modulation is
\begin{equation}
 M_t=\mathbf{1}[\sigma(H_t)\geq0.5],\qquad
 \widehat{\Delta W}_t
 =\alpha\operatorname{Re}\!\left(
 \mathcal{F}^{-1}_2(\bar S_t\odot M_t)\right).
 \label{eq:ums_readout}
\end{equation}
The gate sets entries predicted as off-support to zero.  This decoding rule is
approximate and does not establish exact reconstruction or absence of
inter-task interactions.

\subsection{Dataset-Level Mahalanobis Routing and Inference}
\label{sec:routing}

The router selects which stored UMS coordinate and task head should be decoded
for a grouped query set.  We apply the Gaussian-memory expert-selection
principle used by LEAR~\cite{yu2025lear}.  At the end of task $t$, the complete
training loader is passed through $f_{\mathrm{route}}$.
If $h_0(x)\in\mathbb{R}^{768}$ is its
class token and $u_t(x)=b_t(h_0(x))\in\mathbb{R}^{500}$, the stored statistics
are
\begin{equation}
 \mu_t=\frac{1}{N_t}\sum_i u_t(x_i^t),\qquad
 \Sigma_t=\operatorname{Cov}\{u_t(x_i^t)\}+10^{-4}I.
 \label{eq:gaussian_memory}
\end{equation}
The diagonal term provides numerical regularization.

The main protocol receives an unlabeled query loader $\mathcal{Q}$ whose
samples are already grouped by source dataset/domain.  To match the evaluator,
write $\mathcal{Q}=\{\mathcal{B}_1,\ldots,\mathcal{B}_J\}$ for its mini-batches.
For candidate task $k$, the evaluator first computes the mean of the square-root
Mahalanobis distances within each batch and then accumulates these batch means
with batch-size weights:
\begin{align}
 d_k(x)&=\sqrt{(u_k(x)-\mu_k)^{\mathsf T}\Sigma_k^{-1}
 (u_k(x)-\mu_k)},\notag\\
 D_{k,j}&=\frac{1}{|\mathcal{B}_j|}\sum_{x\in\mathcal{B}_j}d_k(x),\qquad
 \bar D_k(\mathcal{Q})=
 \frac{\sum_{j=1}^{J}|\mathcal{B}_j|D_{k,j}}
      {\sum_{j=1}^{J}|\mathcal{B}_j|},\notag\\
 k^\star&=\arg\min_k\bar D_k(\mathcal{Q}).
 \label{eq:dataset_router}
\end{align}
Here labels are not used in the distance calculation.  The final $k^\star$
after all batches is assigned to the complete loader; the evaluator does not
route each image independently.  The 24 UMS modules decode $z_{k^\star}$ using
Eq.~\eqref{eq:ums_readout}, and prediction is
\begin{equation}
 \widehat y=c_{k^\star}\!\left(
 b_{k^\star}\bigl(f_{\mathrm{loc}}(x;\widehat{\Delta W}_{k^\star})\bigr)
 \right).
 \label{eq:prediction}
\end{equation}

\subsection{Storage and Computational Characteristics}
\label{sec:storage}

One UMS module stores the fixed shared state
$P,Q,W_L,W_R,W_{\mathrm{rot}},\Gamma$, amounting to
$4dR_{\max}+2d^2$ scalars; with $d=R_{\max}=768$, the 24 modules contain
$84{,}934{,}656$ such scalars.  The UMS code itself grows by
$24R_{\max}=18{,}432$ scalars per task, but the complete task-dependent state
also includes the bottleneck, classifier, and 500-dimensional Gaussian
statistics.  Moreover, the present checkpoint retains the preallocated raw
coefficient tensor.  Hence $n_f/R_{\max}\approx7.68$ characterizes only the
code-to-spectrum ratio for one module, not end-to-end model storage.  Training
and consolidation additionally invoke inverse FFTs, dense SVD/QR operations,
Cayley solves, and support decoding.  Total checkpoint size, consolidation
time, and inference latency are therefore measured empirically in the
experiments rather than inferred from the dimension of $z_t$ alone.
.
% Required packages in the main file: amsmath, amssymb, graphicx.

\section{Method}
\label{sec:method}

\subsection{Problem Setup and Framework Overview}
\label{sec:setup}

We consider exemplar-free multi-domain continual learning over a sequence of
$T$ datasets
$\mathcal{D}_{1:T}=\{\mathcal{D}_1,\ldots,\mathcal{D}_T\}$, where
$\mathcal{D}_t=\{(x_i^t,y_i^t)\}_{i=1}^{N_t}$ has a task-specific label
space $\mathcal{Y}_t$.  At step $t$, only $\mathcal{D}_t$ is available for
optimization; images from $\mathcal{D}_{<t}$ are neither stored nor replayed.
The evaluated sequence contains ten heterogeneous visual datasets and therefore
induces both pronounced domain shifts and different prediction heads.

The implementation instantiates two separately pretrained and frozen ViT-B/16
networks~\cite{dosovitskiy2021vit}.  The local backbone
$f_{\mathrm{loc}}$ is used for classification and receives additive modulation
in selected attention projections.  The routing backbone $f_{\mathrm{route}}$
remains unmodulated and provides a class-token feature for task selection.  For
each task $t$, a linear bottleneck $b_t:\mathbb{R}^{768}\rightarrow
\mathbb{R}^{500}$ and a classifier
$c_t:\mathbb{R}^{500}\rightarrow\mathbb{R}^{|\mathcal{Y}_t|}$ are learned and
then frozen.

FOCAL-UMS follows an encode--decode process.  Task-specific Fourier
coefficients first generate additive attention modulations.  After training,
each modulation is encoded as a fixed-length task coordinate in an incrementally
constructed \emph{Unified Modulation Subspace} (UMS); at inference, that
coordinate is decoded into a modulation matrix.  Subspace-aware training and
post-task refinement connect these two stages by exposing the task parameters
to the candidate representation during learning and by fitting the shared
rotation and support decoder after encoding.  Sparse Fourier parameterization
and Mahalanobis routing are adopted mechanisms.  The methodological focus of
FOCAL-UMS is the integration of UMS encoding and decoding into sequential
learning; the refinement objective is one component of this process rather than
a stand-alone collection of loss functions.

\begin{figure}[t]
  \centering
  \includegraphics[width=\linewidth]{\detokenize{FOCAL-UMS pipeline with post-task refinement.png}}
  \caption{FOCAL-UMS.  Task-specific Fourier coefficients are optimized through
  a candidate UMS map, encoded after training, and refined before decoding.  At
  inference, one task is selected for the complete grouped query set and its
  UMS code is decoded into attention modulations.  The anchor regularizes old
  rotation directions; it does not imply invariant old-task modulations.}
  \label{fig:focal_ums}
\end{figure}

\subsection{Task-Specific Frequency Modulation}
\label{sec:frequency}

We adapt sparse Fourier-domain fine-tuning~\cite{gao2024fourierft} to allocate
separate trainable coordinates to successive tasks.  Let $d=768$ and let $\pi$
be a deterministic, seeded permutation of the $d^2$ entries of a square
spectrum.  The first $Tn_f$ entries are partitioned into pairwise-disjoint
supports $\Omega_1,\ldots,\Omega_T$, each of cardinality $n_f$.  For task $t$,
the vector $\boldsymbol{\theta}_t\in\mathbb{R}^{n_f}$ is scattered onto
$\Omega_t$ to form
\begin{equation}
 S_t[u_m,v_m]=\theta_{t,m}\quad ((u_m,v_m)\in\Omega_t),
 \qquad S_t[u,v]=0\quad ((u,v)\notin\Omega_t).
 \label{eq:sparse_spectrum}
\end{equation}
The corresponding additive weight modulation is
\begin{equation}
 \Delta W_t=\alpha\,\operatorname{Re}\!\left(\mathcal{F}^{-1}_2(S_t)\right),
 \label{eq:ifft_modulation}
\end{equation}
where $\mathcal{F}^{-1}_2$ follows the default normalization of
\texttt{torch.fft.ifft2}.  In the primary configuration, $n_f=5{,}898$ and
$\alpha=300$.  The model modulates the query, key, value, and output projections
in blocks $6$--$11$, yielding 24 independent FocalExpert/UMS modules.  Their
pretrained projection weights remain frozen.

Disjoint supports prevent two tasks from directly updating the same allocated
raw spectral coordinates.  This statement is deliberately local to the raw
parameterization: taking the real part in Eq.~\eqref{eq:ifft_modulation}, and
later applying UMS reconstruction and support gating, does not establish exact
orthogonality of the resulting spatial updates or guarantee the absence of
forgetting.

\subsection{UMS-in-the-Loop Training}
\label{sec:inloop}

Each modulated projection owns a UMS with column bases
$P,Q\in\mathbb{R}^{d\times R_{\max}}$, a current active rank $r_t$, and a
shared rotation $O\in\mathbb{R}^{d\times d}$.  UMS-in-the-loop is designed to
reduce, rather than eliminate, the mismatch between unconstrained task training
and subsequent subspace encoding.  At task start and at periodic refreshes, the
current spectrum is transformed to the canonical frame,
\begin{equation}
 T_t=O^{\mathsf T}S_tO.
 \label{eq:canonical}
\end{equation}
The implementation removes the component explained by the active left basis,
computes a truncated SVD of the residual, residualizes the resulting right
directions against $Q$, and orthonormalizes both sides by QR.  Appending these
directions to the active columns gives candidate bases
$P_{\mathrm{c}},Q_{\mathrm{c}}\in\mathbb{R}^{d\times k}$, capped by
$R_{\max}$.

Let $A=OP_{\mathrm{c}}$ and $B=OQ_{\mathrm{c}}$.  Restricting the induced
bilinear representation to the current support gives
\begin{equation}
 G_t[m,j]=A[u_m,j]B[v_m,j],
 \qquad G_t\in\mathbb{R}^{n_f\times k}.
 \label{eq:G}
\end{equation}
The coefficients used in the forward pass are
\begin{equation}
 \widetilde{\boldsymbol{\theta}}_t
 =G_t\bigl(G_t^{\mathsf T}\boldsymbol{\theta}_t\bigr),
 \label{eq:inloop_map}
\end{equation}
and are scattered onto $\Omega_t$ before Eq.~\eqref{eq:ifft_modulation}.  The
matrix $G_t$ is cached without gradients and is refreshed at task start and
periodically during training.  Because its support-restricted columns are not
explicitly orthonormal and the code does not apply
$(G_t^{\mathsf T}G_t)^{-1}$, Eq.~\eqref{eq:inloop_map} is a
candidate-subspace-induced linear map, not an orthogonal projector.  Cross
entropy updates the current raw coefficients, bottleneck, and classifier.

\subsection{Post-Task UMS Encoding}
\label{sec:encoding}

The purpose of post-task encoding is to express each learned modulation through
shared bases and a task-specific coordinate that can later be decoded.  At task
completion, the current implementation rebuilds $S_t$ from the
\emph{raw} slice $\boldsymbol{\theta}_t$ and encodes this matrix into UMS.  It
does not fold the final in-loop value
$\widetilde{\boldsymbol{\theta}}_t$ from Eq.~\eqref{eq:inloop_map}; thus the
in-loop mechanism biases learning toward the candidate representation but does
not make training and end-task encoding identical.

Let $r_{t-1}$ be the active rank before encoding task $t$, with $r_0=0$.
For the previously active bases $P_{t-1}=P_{:,1:r_{t-1}}$ and
$Q_{t-1}=Q_{:,1:r_{t-1}}$, the left residual is
\begin{equation}
 U_t=T_t-P_{t-1}(P_{t-1}^{\mathsf T}T_t).
 \label{eq:left_residual}
\end{equation}
A truncated SVD
$U_t\approx U_s\operatorname{diag}(\boldsymbol{\sigma})V_s^{\mathsf T}$
provides paired residual directions.  The right directions are additionally
residualized against $Q_{t-1}$; QR-orthonormalized directions are appended to $P$
and $Q$ when
\begin{equation}
 \mathcal{E}_t=\|U_s\operatorname{diag}(\boldsymbol{\sigma})\|_F^2>\tau
 \label{eq:residual_energy}
\end{equation}
and capacity remains.  In the evaluated configuration,
$R_{\max}=768$, the truncated-SVD budget is 75, and $\tau=10^{-3}$.  With the
updated active rank $r_t$ and bases $P_r=P_{:,1:r_t}$ and
$Q_r=Q_{:,1:r_t}$, task $t$ is encoded as the zero-padded diagonal core
\begin{equation}
 z_t=\operatorname{pad}_{R_{\max}}\!\left[
 \operatorname{diag}\!\left(P_r^{\mathsf T}T_tQ_r\right)\right].
 \label{eq:task_code}
\end{equation}
This is a rank-budgeted diagonal approximation, not an exact factorization of
the complete core.  UMS codes are used for inference, although the present
training/checkpoint implementation also retains the preallocated raw
coefficient tensor.

\subsection{UMS Refinement and Decoding}
\label{sec:multiobjective}

The purpose of refinement is to reduce the discrepancy between the sparse task
spectrum and the spectrum decoded from the diagonal UMS coordinate.  A single
reconstruction term is insufficient for this purpose: it does not identify the
sparse support, expose off-diagonal energy omitted by $z_t$, or control changes
to previously used rotation directions.  We therefore optimize the rotation
and support decoder with a structured objective derived from these encoding
errors.

The shared orientation is parameterized by the Cayley transform
\begin{equation}
 A_{\mathrm{rot}}=W_{\mathrm{rot}}-W_{\mathrm{rot}}^{\mathsf T},\qquad
 O=\left(I-\tfrac{1}{2}A_{\mathrm{rot}}\right)^{-1}
   \left(I+\tfrac{1}{2}A_{\mathrm{rot}}\right),
 \label{eq:cayley}
\end{equation}
which is orthogonal in exact arithmetic whenever the solve is well-defined.
For the current code $z_t$, define the ungated spectral reconstruction and the
support logits as
\begin{align}
 \bar S_t &= OP\operatorname{diag}(z_t)Q^{\mathsf T}O^{\mathsf T},
 \label{eq:ungated_reconstruction}\\
 H_t &= W_L\operatorname{diag}(z_t)W_R^{\mathsf T}+\Gamma.
 \label{eq:support_logits}
\end{align}
Here $W_L,W_R\in\mathbb{R}^{d\times R_{\max}}$ and
$\Gamma\in\mathbb{R}^{d\times d}$ are shared support-decoder parameters.  Let
$\mathcal{I}_t=\{(u,v):S_t[u,v]\neq0\}$,
$Y_t[u,v]=\mathbf{1}[(u,v)\in\mathcal{I}_t]$, and
$C_t=P_r^{\mathsf T}(O^{\mathsf T}S_tO)Q_r$.  For nonempty $\mathcal{I}_t$,
let $m_t$ be the mean of $\bar S_t$ on $\mathcal{I}_t$ and let
$s_t=\sqrt{\operatorname{Var}_{\mathcal{I}_t}(\bar S_t)+10^{-8}}$,
where the variance uses the population denominator.  The code $z_t$ is held
fixed during refinement.  For $t>1$, let $O^{-},P^{-},Q^{-}$ denote the
pre-refinement rotation and previously active basis columns, with
$r_{t-1}$ columns.  The effective refinement objective is
\begin{align}
 \mathcal{L}_{\mathrm{UMS}}
 &=\lambda_{\mathrm{val}}
   \frac{\sum_{(u,v)\in\mathcal{I}_t}(S_t[u,v]-\bar S_t[u,v])^2}
        {|\mathcal{I}_t|}
 +\lambda_{\mathrm{mask}}\operatorname{BCEWithLogits}(H_t,Y_t;w_+)
 \notag\\
 &\quad+\lambda_{\mathrm{mom}}\bigl(m_t^2+(s_t-1)^2\bigr)
 +\lambda_{\mathrm{off}}
   \frac{\|C_t-\operatorname{diag}(\operatorname{diag}C_t)\|_F^2}
        {\max\{r_t(r_t-1),1\}}
 \notag\\
 &\quad+\lambda_{\mathrm{drift}}
   \frac{\|(O-O^{-})P^{-}\|_F^2+\|(O-O^{-})Q^{-}\|_F^2}
        {2r_{t-1}},
 \label{eq:loss_total}
\end{align}
where $w_+=(d^2-n_f)/n_f$ compensates for support imbalance.  The first two
terms penalize value and support-mask errors, respectively; the moment term
regularizes the scale of decoded occupied entries; the off-diagonal term
penalizes canonical-core energy omitted by the diagonal code; and the drift
term penalizes rotation changes along previously active basis directions.  The
value and moment terms are zero when $\mathcal{I}_t$ is empty; the drift term
is zero when there are no previously active columns.  The implemented weights are
$(\lambda_{\mathrm{val}},\lambda_{\mathrm{mask}},\lambda_{\mathrm{mom}},
\lambda_{\mathrm{off}},\lambda_{\mathrm{drift}})
=(1.0,0.1,0.01,0.1,2.0)$.  Refinement uses 15 Adam steps with learning rates
$10^{-4}$ for $W_{\mathrm{rot}}$ and $10^{-3}$ for the support decoder, after
which the current task is re-encoded without further basis expansion.

The implementation also forms a nominal drift penalty for the previously
active columns of $W_L$ and $W_R$.  Its reference tensors are detached views,
however, rather than independent copies: they share storage with the current
parameters and change when Adam updates them.  Consequently, their difference
from the current columns and the resulting penalty and gradient are zero.
Equation~\eqref{eq:loss_total} therefore states the effective drift term.  The
anchor does not constrain $\Gamma$, re-encode all previous tasks, or establish
invariant old-task readouts.  Numerical orthogonality of the Cayley rotation
can be assessed using the saved diagnostics.

At readout, the support decoder produces a hard gate and the task modulation is
\begin{equation}
 M_t=\mathbf{1}[\sigma(H_t)\geq0.5],\qquad
 \widehat{\Delta W}_t
 =\alpha\operatorname{Re}\!\left(
 \mathcal{F}^{-1}_2(\bar S_t\odot M_t)\right).
 \label{eq:ums_readout}
\end{equation}
The gate sets entries predicted as off-support to zero.  This decoding rule is
approximate and does not establish exact reconstruction or absence of
inter-task interactions.

\subsection{Dataset-Level Mahalanobis Routing and Inference}
\label{sec:routing}

The router selects which stored UMS coordinate and task head should be decoded
for a grouped query set.  We apply the Gaussian-memory expert-selection
principle used by LEAR~\cite{yu2025lear}.  At the end of task $t$, the complete
training loader is passed through $f_{\mathrm{route}}$.
If $h_0(x)\in\mathbb{R}^{768}$ is its
class token and $u_t(x)=b_t(h_0(x))\in\mathbb{R}^{500}$, the stored statistics
are
\begin{equation}
 \mu_t=\frac{1}{N_t}\sum_i u_t(x_i^t),\qquad
 \Sigma_t=\operatorname{Cov}\{u_t(x_i^t)\}+10^{-4}I.
 \label{eq:gaussian_memory}
\end{equation}
The diagonal term provides numerical regularization.

The main protocol receives an unlabeled query loader $\mathcal{Q}$ whose
samples are already grouped by source dataset/domain.  To match the evaluator,
write $\mathcal{Q}=\{\mathcal{B}_1,\ldots,\mathcal{B}_J\}$ for its mini-batches.
For candidate task $k$, the evaluator first computes the mean of the square-root
Mahalanobis distances within each batch and then accumulates these batch means
with batch-size weights:
\begin{align}
 d_k(x)&=\sqrt{(u_k(x)-\mu_k)^{\mathsf T}\Sigma_k^{-1}
 (u_k(x)-\mu_k)},\notag\\
 D_{k,j}&=\frac{1}{|\mathcal{B}_j|}\sum_{x\in\mathcal{B}_j}d_k(x),\qquad
 \bar D_k(\mathcal{Q})=
 \frac{\sum_{j=1}^{J}|\mathcal{B}_j|D_{k,j}}
      {\sum_{j=1}^{J}|\mathcal{B}_j|},\notag\\
 k^\star&=\arg\min_k\bar D_k(\mathcal{Q}).
 \label{eq:dataset_router}
\end{align}
Here labels are not used in the distance calculation.  The final $k^\star$
after all batches is assigned to the complete loader; the evaluator does not
route each image independently.  The 24 UMS modules decode $z_{k^\star}$ using
Eq.~\eqref{eq:ums_readout}, and prediction is
\begin{equation}
 \widehat y=c_{k^\star}\!\left(
 b_{k^\star}\bigl(f_{\mathrm{loc}}(x;\widehat{\Delta W}_{k^\star})\bigr)
 \right).
 \label{eq:prediction}
\end{equation}

\subsection{Storage and Computational Characteristics}
\label{sec:storage}

One UMS module stores the fixed shared state
$P,Q,W_L,W_R,W_{\mathrm{rot}},\Gamma$, amounting to
$4dR_{\max}+2d^2$ scalars; with $d=R_{\max}=768$, the 24 modules contain
$84{,}934{,}656$ such scalars.  The UMS code itself grows by
$24R_{\max}=18{,}432$ scalars per task, but the complete task-dependent state
also includes the bottleneck, classifier, and 500-dimensional Gaussian
statistics.  Moreover, the present checkpoint retains the preallocated raw
coefficient tensor.  Hence $n_f/R_{\max}\approx7.68$ characterizes only the
code-to-spectrum ratio for one module, not end-to-end model storage.  Training
and consolidation additionally invoke inverse FFTs, dense SVD/QR operations,
Cayley solves, and support decoding.  Total checkpoint size, consolidation
time, and inference latency are therefore measured empirically in the
experiments rather than inferred from the dimension of $z_t$ alone.
.
% Required packages in the main file: amsmath, amssymb, graphicx.

\section{Method}
\label{sec:method}

\subsection{Problem Setup and Framework Overview}
\label{sec:setup}

We consider exemplar-free multi-domain continual learning over a sequence of
$T$ datasets
$\mathcal{D}_{1:T}=\{\mathcal{D}_1,\ldots,\mathcal{D}_T\}$, where
$\mathcal{D}_t=\{(x_i^t,y_i^t)\}_{i=1}^{N_t}$ has a task-specific label
space $\mathcal{Y}_t$.  At step $t$, only $\mathcal{D}_t$ is available for
optimization; images from $\mathcal{D}_{<t}$ are neither stored nor replayed.
The evaluated sequence contains ten heterogeneous visual datasets and therefore
induces both pronounced domain shifts and different prediction heads.

The implementation instantiates two separately pretrained and frozen ViT-B/16
networks~\cite{dosovitskiy2021vit}.  The local backbone
$f_{\mathrm{loc}}$ is used for classification and receives additive modulation
in selected attention projections.  The routing backbone $f_{\mathrm{route}}$
remains unmodulated and provides a class-token feature for task selection.  For
each task $t$, a linear bottleneck $b_t:\mathbb{R}^{768}\rightarrow
\mathbb{R}^{500}$ and a classifier
$c_t:\mathbb{R}^{500}\rightarrow\mathbb{R}^{|\mathcal{Y}_t|}$ are learned and
then frozen.

FOCAL-UMS follows an encode--decode process.  Task-specific Fourier
coefficients first generate additive attention modulations.  After training,
each modulation is encoded as a fixed-length task coordinate in an incrementally
constructed \emph{Unified Modulation Subspace} (UMS); at inference, that
coordinate is decoded into a modulation matrix.  Subspace-aware training and
post-task refinement connect these two stages by exposing the task parameters
to the candidate representation during learning and by fitting the shared
rotation and support decoder after encoding.  Sparse Fourier parameterization
and Mahalanobis routing are adopted mechanisms.  The methodological focus of
FOCAL-UMS is the integration of UMS encoding and decoding into sequential
learning; the refinement objective is one component of this process rather than
a stand-alone collection of loss functions.

\begin{figure}[t]
  \centering
  \includegraphics[width=\linewidth]{\detokenize{FOCAL-UMS pipeline with post-task refinement.png}}
  \caption{FOCAL-UMS.  Task-specific Fourier coefficients are optimized through
  a candidate UMS map, encoded after training, and refined before decoding.  At
  inference, one task is selected for the complete grouped query set and its
  UMS code is decoded into attention modulations.  The anchor regularizes old
  rotation directions; it does not imply invariant old-task modulations.}
  \label{fig:focal_ums}
\end{figure}

\subsection{Task-Specific Frequency Modulation}
\label{sec:frequency}

We adapt sparse Fourier-domain fine-tuning~\cite{gao2024fourierft} to allocate
separate trainable coordinates to successive tasks.  Let $d=768$ and let $\pi$
be a deterministic, seeded permutation of the $d^2$ entries of a square
spectrum.  The first $Tn_f$ entries are partitioned into pairwise-disjoint
supports $\Omega_1,\ldots,\Omega_T$, each of cardinality $n_f$.  For task $t$,
the vector $\boldsymbol{\theta}_t\in\mathbb{R}^{n_f}$ is scattered onto
$\Omega_t$ to form
\begin{equation}
 S_t[u_m,v_m]=\theta_{t,m}\quad ((u_m,v_m)\in\Omega_t),
 \qquad S_t[u,v]=0\quad ((u,v)\notin\Omega_t).
 \label{eq:sparse_spectrum}
\end{equation}
The corresponding additive weight modulation is
\begin{equation}
 \Delta W_t=\alpha\,\operatorname{Re}\!\left(\mathcal{F}^{-1}_2(S_t)\right),
 \label{eq:ifft_modulation}
\end{equation}
where $\mathcal{F}^{-1}_2$ follows the default normalization of
\texttt{torch.fft.ifft2}.  In the primary configuration, $n_f=5{,}898$ and
$\alpha=300$.  The model modulates the query, key, value, and output projections
in blocks $6$--$11$, yielding 24 independent FocalExpert/UMS modules.  Their
pretrained projection weights remain frozen.

Disjoint supports prevent two tasks from directly updating the same allocated
raw spectral coordinates.  This statement is deliberately local to the raw
parameterization: taking the real part in Eq.~\eqref{eq:ifft_modulation}, and
later applying UMS reconstruction and support gating, does not establish exact
orthogonality of the resulting spatial updates or guarantee the absence of
forgetting.

\subsection{UMS-in-the-Loop Training}
\label{sec:inloop}

Each modulated projection owns a UMS with column bases
$P,Q\in\mathbb{R}^{d\times R_{\max}}$, a current active rank $r_t$, and a
shared rotation $O\in\mathbb{R}^{d\times d}$.  UMS-in-the-loop is designed to
reduce, rather than eliminate, the mismatch between unconstrained task training
and subsequent subspace encoding.  At task start and at periodic refreshes, the
current spectrum is transformed to the canonical frame,
\begin{equation}
 T_t=O^{\mathsf T}S_tO.
 \label{eq:canonical}
\end{equation}
The implementation removes the component explained by the active left basis,
computes a truncated SVD of the residual, residualizes the resulting right
directions against $Q$, and orthonormalizes both sides by QR.  Appending these
directions to the active columns gives candidate bases
$P_{\mathrm{c}},Q_{\mathrm{c}}\in\mathbb{R}^{d\times k}$, capped by
$R_{\max}$.

Let $A=OP_{\mathrm{c}}$ and $B=OQ_{\mathrm{c}}$.  Restricting the induced
bilinear representation to the current support gives
\begin{equation}
 G_t[m,j]=A[u_m,j]B[v_m,j],
 \qquad G_t\in\mathbb{R}^{n_f\times k}.
 \label{eq:G}
\end{equation}
The coefficients used in the forward pass are
\begin{equation}
 \widetilde{\boldsymbol{\theta}}_t
 =G_t\bigl(G_t^{\mathsf T}\boldsymbol{\theta}_t\bigr),
 \label{eq:inloop_map}
\end{equation}
and are scattered onto $\Omega_t$ before Eq.~\eqref{eq:ifft_modulation}.  The
matrix $G_t$ is cached without gradients and is refreshed at task start and
periodically during training.  Because its support-restricted columns are not
explicitly orthonormal and the code does not apply
$(G_t^{\mathsf T}G_t)^{-1}$, Eq.~\eqref{eq:inloop_map} is a
candidate-subspace-induced linear map, not an orthogonal projector.  Cross
entropy updates the current raw coefficients, bottleneck, and classifier.

\subsection{Post-Task UMS Encoding}
\label{sec:encoding}

The purpose of post-task encoding is to express each learned modulation through
shared bases and a task-specific coordinate that can later be decoded.  At task
completion, the current implementation rebuilds $S_t$ from the
\emph{raw} slice $\boldsymbol{\theta}_t$ and encodes this matrix into UMS.  It
does not fold the final in-loop value
$\widetilde{\boldsymbol{\theta}}_t$ from Eq.~\eqref{eq:inloop_map}; thus the
in-loop mechanism biases learning toward the candidate representation but does
not make training and end-task encoding identical.

Let $r_{t-1}$ be the active rank before encoding task $t$, with $r_0=0$.
For the previously active bases $P_{t-1}=P_{:,1:r_{t-1}}$ and
$Q_{t-1}=Q_{:,1:r_{t-1}}$, the left residual is
\begin{equation}
 U_t=T_t-P_{t-1}(P_{t-1}^{\mathsf T}T_t).
 \label{eq:left_residual}
\end{equation}
A truncated SVD
$U_t\approx U_s\operatorname{diag}(\boldsymbol{\sigma})V_s^{\mathsf T}$
provides paired residual directions.  The right directions are additionally
residualized against $Q_{t-1}$; QR-orthonormalized directions are appended to $P$
and $Q$ when
\begin{equation}
 \mathcal{E}_t=\|U_s\operatorname{diag}(\boldsymbol{\sigma})\|_F^2>\tau
 \label{eq:residual_energy}
\end{equation}
and capacity remains.  In the evaluated configuration,
$R_{\max}=768$, the truncated-SVD budget is 75, and $\tau=10^{-3}$.  With the
updated active rank $r_t$ and bases $P_r=P_{:,1:r_t}$ and
$Q_r=Q_{:,1:r_t}$, task $t$ is encoded as the zero-padded diagonal core
\begin{equation}
 z_t=\operatorname{pad}_{R_{\max}}\!\left[
 \operatorname{diag}\!\left(P_r^{\mathsf T}T_tQ_r\right)\right].
 \label{eq:task_code}
\end{equation}
This is a rank-budgeted diagonal approximation, not an exact factorization of
the complete core.  UMS codes are used for inference, although the present
training/checkpoint implementation also retains the preallocated raw
coefficient tensor.

\subsection{UMS Refinement and Decoding}
\label{sec:multiobjective}

The purpose of refinement is to reduce the discrepancy between the sparse task
spectrum and the spectrum decoded from the diagonal UMS coordinate.  A single
reconstruction term is insufficient for this purpose: it does not identify the
sparse support, expose off-diagonal energy omitted by $z_t$, or control changes
to previously used rotation directions.  We therefore optimize the rotation
and support decoder with a structured objective derived from these encoding
errors.

The shared orientation is parameterized by the Cayley transform
\begin{equation}
 A_{\mathrm{rot}}=W_{\mathrm{rot}}-W_{\mathrm{rot}}^{\mathsf T},\qquad
 O=\left(I-\tfrac{1}{2}A_{\mathrm{rot}}\right)^{-1}
   \left(I+\tfrac{1}{2}A_{\mathrm{rot}}\right),
 \label{eq:cayley}
\end{equation}
which is orthogonal in exact arithmetic whenever the solve is well-defined.
For the current code $z_t$, define the ungated spectral reconstruction and the
support logits as
\begin{align}
 \bar S_t &= OP\operatorname{diag}(z_t)Q^{\mathsf T}O^{\mathsf T},
 \label{eq:ungated_reconstruction}\\
 H_t &= W_L\operatorname{diag}(z_t)W_R^{\mathsf T}+\Gamma.
 \label{eq:support_logits}
\end{align}
Here $W_L,W_R\in\mathbb{R}^{d\times R_{\max}}$ and
$\Gamma\in\mathbb{R}^{d\times d}$ are shared support-decoder parameters.  Let
$\mathcal{I}_t=\{(u,v):S_t[u,v]\neq0\}$,
$Y_t[u,v]=\mathbf{1}[(u,v)\in\mathcal{I}_t]$, and
$C_t=P_r^{\mathsf T}(O^{\mathsf T}S_tO)Q_r$.  For nonempty $\mathcal{I}_t$,
let $m_t$ be the mean of $\bar S_t$ on $\mathcal{I}_t$ and let
$s_t=\sqrt{\operatorname{Var}_{\mathcal{I}_t}(\bar S_t)+10^{-8}}$,
where the variance uses the population denominator.  The code $z_t$ is held
fixed during refinement.  For $t>1$, let $O^{-},P^{-},Q^{-}$ denote the
pre-refinement rotation and previously active basis columns, with
$r_{t-1}$ columns.  The effective refinement objective is
\begin{align}
 \mathcal{L}_{\mathrm{UMS}}
 &=\lambda_{\mathrm{val}}
   \frac{\sum_{(u,v)\in\mathcal{I}_t}(S_t[u,v]-\bar S_t[u,v])^2}
        {|\mathcal{I}_t|}
 +\lambda_{\mathrm{mask}}\operatorname{BCEWithLogits}(H_t,Y_t;w_+)
 \notag\\
 &\quad+\lambda_{\mathrm{mom}}\bigl(m_t^2+(s_t-1)^2\bigr)
 +\lambda_{\mathrm{off}}
   \frac{\|C_t-\operatorname{diag}(\operatorname{diag}C_t)\|_F^2}
        {\max\{r_t(r_t-1),1\}}
 \notag\\
 &\quad+\lambda_{\mathrm{drift}}
   \frac{\|(O-O^{-})P^{-}\|_F^2+\|(O-O^{-})Q^{-}\|_F^2}
        {2r_{t-1}},
 \label{eq:loss_total}
\end{align}
where $w_+=(d^2-n_f)/n_f$ compensates for support imbalance.  The first two
terms penalize value and support-mask errors, respectively; the moment term
regularizes the scale of decoded occupied entries; the off-diagonal term
penalizes canonical-core energy omitted by the diagonal code; and the drift
term penalizes rotation changes along previously active basis directions.  The
value and moment terms are zero when $\mathcal{I}_t$ is empty; the drift term
is zero when there are no previously active columns.  The implemented weights are
$(\lambda_{\mathrm{val}},\lambda_{\mathrm{mask}},\lambda_{\mathrm{mom}},
\lambda_{\mathrm{off}},\lambda_{\mathrm{drift}})
=(1.0,0.1,0.01,0.1,2.0)$.  Refinement uses 15 Adam steps with learning rates
$10^{-4}$ for $W_{\mathrm{rot}}$ and $10^{-3}$ for the support decoder, after
which the current task is re-encoded without further basis expansion.

The implementation also forms a nominal drift penalty for the previously
active columns of $W_L$ and $W_R$.  Its reference tensors are detached views,
however, rather than independent copies: they share storage with the current
parameters and change when Adam updates them.  Consequently, their difference
from the current columns and the resulting penalty and gradient are zero.
Equation~\eqref{eq:loss_total} therefore states the effective drift term.  The
anchor does not constrain $\Gamma$, re-encode all previous tasks, or establish
invariant old-task readouts.  Numerical orthogonality of the Cayley rotation
can be assessed using the saved diagnostics.

At readout, the support decoder produces a hard gate and the task modulation is
\begin{equation}
 M_t=\mathbf{1}[\sigma(H_t)\geq0.5],\qquad
 \widehat{\Delta W}_t
 =\alpha\operatorname{Re}\!\left(
 \mathcal{F}^{-1}_2(\bar S_t\odot M_t)\right).
 \label{eq:ums_readout}
\end{equation}
The gate sets entries predicted as off-support to zero.  This decoding rule is
approximate and does not establish exact reconstruction or absence of
inter-task interactions.

\subsection{Dataset-Level Mahalanobis Routing and Inference}
\label{sec:routing}

The router selects which stored UMS coordinate and task head should be decoded
for a grouped query set.  We apply the Gaussian-memory expert-selection
principle used by LEAR~\cite{yu2025lear}.  At the end of task $t$, the complete
training loader is passed through $f_{\mathrm{route}}$.
If $h_0(x)\in\mathbb{R}^{768}$ is its
class token and $u_t(x)=b_t(h_0(x))\in\mathbb{R}^{500}$, the stored statistics
are
\begin{equation}
 \mu_t=\frac{1}{N_t}\sum_i u_t(x_i^t),\qquad
 \Sigma_t=\operatorname{Cov}\{u_t(x_i^t)\}+10^{-4}I.
 \label{eq:gaussian_memory}
\end{equation}
The diagonal term provides numerical regularization.

The main protocol receives an unlabeled query loader $\mathcal{Q}$ whose
samples are already grouped by source dataset/domain.  To match the evaluator,
write $\mathcal{Q}=\{\mathcal{B}_1,\ldots,\mathcal{B}_J\}$ for its mini-batches.
For candidate task $k$, the evaluator first computes the mean of the square-root
Mahalanobis distances within each batch and then accumulates these batch means
with batch-size weights:
\begin{align}
 d_k(x)&=\sqrt{(u_k(x)-\mu_k)^{\mathsf T}\Sigma_k^{-1}
 (u_k(x)-\mu_k)},\notag\\
 D_{k,j}&=\frac{1}{|\mathcal{B}_j|}\sum_{x\in\mathcal{B}_j}d_k(x),\qquad
 \bar D_k(\mathcal{Q})=
 \frac{\sum_{j=1}^{J}|\mathcal{B}_j|D_{k,j}}
      {\sum_{j=1}^{J}|\mathcal{B}_j|},\notag\\
 k^\star&=\arg\min_k\bar D_k(\mathcal{Q}).
 \label{eq:dataset_router}
\end{align}
Here labels are not used in the distance calculation.  The final $k^\star$
after all batches is assigned to the complete loader; the evaluator does not
route each image independently.  The 24 UMS modules decode $z_{k^\star}$ using
Eq.~\eqref{eq:ums_readout}, and prediction is
\begin{equation}
 \widehat y=c_{k^\star}\!\left(
 b_{k^\star}\bigl(f_{\mathrm{loc}}(x;\widehat{\Delta W}_{k^\star})\bigr)
 \right).
 \label{eq:prediction}
\end{equation}

\subsection{Storage and Computational Characteristics}
\label{sec:storage}

One UMS module stores the fixed shared state
$P,Q,W_L,W_R,W_{\mathrm{rot}},\Gamma$, amounting to
$4dR_{\max}+2d^2$ scalars; with $d=R_{\max}=768$, the 24 modules contain
$84{,}934{,}656$ such scalars.  The UMS code itself grows by
$24R_{\max}=18{,}432$ scalars per task, but the complete task-dependent state
also includes the bottleneck, classifier, and 500-dimensional Gaussian
statistics.  Moreover, the present checkpoint retains the preallocated raw
coefficient tensor.  Hence $n_f/R_{\max}\approx7.68$ characterizes only the
code-to-spectrum ratio for one module, not end-to-end model storage.  Training
and consolidation additionally invoke inverse FFTs, dense SVD/QR operations,
Cayley solves, and support decoding.  Total checkpoint size, consolidation
time, and inference latency are therefore measured empirically in the
experiments rather than inferred from the dimension of $z_t$ alone.
.
% Required packages in the main file: amsmath, amssymb, graphicx.

\section{Method}
\label{sec:method}

\subsection{Problem Setup and Framework Overview}
\label{sec:setup}

We consider exemplar-free multi-domain continual learning over a sequence of
$T$ datasets
$\mathcal{D}_{1:T}=\{\mathcal{D}_1,\ldots,\mathcal{D}_T\}$, where
$\mathcal{D}_t=\{(x_i^t,y_i^t)\}_{i=1}^{N_t}$ has a task-specific label
space $\mathcal{Y}_t$.  At step $t$, only $\mathcal{D}_t$ is available for
optimization; images from $\mathcal{D}_{<t}$ are neither stored nor replayed.
The evaluated sequence contains ten heterogeneous visual datasets and therefore
induces both pronounced domain shifts and different prediction heads.

The implementation instantiates two separately pretrained and frozen ViT-B/16
networks~\cite{dosovitskiy2021vit}.  The local backbone
$f_{\mathrm{loc}}$ is used for classification and receives additive modulation
in selected attention projections.  The routing backbone $f_{\mathrm{route}}$
remains unmodulated and provides a class-token feature for task selection.  For
each task $t$, a linear bottleneck $b_t:\mathbb{R}^{768}\rightarrow
\mathbb{R}^{500}$ and a classifier
$c_t:\mathbb{R}^{500}\rightarrow\mathbb{R}^{|\mathcal{Y}_t|}$ are learned and
then frozen.

FOCAL-UMS follows an encode--decode process.  Task-specific Fourier
coefficients first generate additive attention modulations.  After training,
each modulation is encoded as a fixed-length task coordinate in an incrementally
constructed \emph{Unified Modulation Subspace} (UMS); at inference, that
coordinate is decoded into a modulation matrix.  Subspace-aware training and
post-task refinement connect these two stages by exposing the task parameters
to the candidate representation during learning and by fitting the shared
rotation and support decoder after encoding.  Sparse Fourier parameterization
and Mahalanobis routing are adopted mechanisms.  The methodological focus of
FOCAL-UMS is the integration of UMS encoding and decoding into sequential
learning; the refinement objective is one component of this process rather than
a stand-alone collection of loss functions.

\begin{figure}[t]
  \centering
  \includegraphics[width=\linewidth]{\detokenize{FOCAL-UMS pipeline with post-task refinement.png}}
  \caption{FOCAL-UMS.  Task-specific Fourier coefficients are optimized through
  a candidate UMS map, encoded after training, and refined before decoding.  At
  inference, one task is selected for the complete grouped query set and its
  UMS code is decoded into attention modulations.  The anchor regularizes old
  rotation directions; it does not imply invariant old-task modulations.}
  \label{fig:focal_ums}
\end{figure}

\subsection{Task-Specific Frequency Modulation}
\label{sec:frequency}

We adapt sparse Fourier-domain fine-tuning~\cite{gao2024fourierft} to allocate
separate trainable coordinates to successive tasks.  Let $d=768$ and let $\pi$
be a deterministic, seeded permutation of the $d^2$ entries of a square
spectrum.  The first $Tn_f$ entries are partitioned into pairwise-disjoint
supports $\Omega_1,\ldots,\Omega_T$, each of cardinality $n_f$.  For task $t$,
the vector $\boldsymbol{\theta}_t\in\mathbb{R}^{n_f}$ is scattered onto
$\Omega_t$ to form
\begin{equation}
 S_t[u_m,v_m]=\theta_{t,m}\quad ((u_m,v_m)\in\Omega_t),
 \qquad S_t[u,v]=0\quad ((u,v)\notin\Omega_t).
 \label{eq:sparse_spectrum}
\end{equation}
The corresponding additive weight modulation is
\begin{equation}
 \Delta W_t=\alpha\,\operatorname{Re}\!\left(\mathcal{F}^{-1}_2(S_t)\right),
 \label{eq:ifft_modulation}
\end{equation}
where $\mathcal{F}^{-1}_2$ follows the default normalization of
\texttt{torch.fft.ifft2}.  In the primary configuration, $n_f=5{,}898$ and
$\alpha=300$.  The model modulates the query, key, value, and output projections
in blocks $6$--$11$, yielding 24 independent FocalExpert/UMS modules.  Their
pretrained projection weights remain frozen.

Disjoint supports prevent two tasks from directly updating the same allocated
raw spectral coordinates.  This statement is deliberately local to the raw
parameterization: taking the real part in Eq.~\eqref{eq:ifft_modulation}, and
later applying UMS reconstruction and support gating, does not establish exact
orthogonality of the resulting spatial updates or guarantee the absence of
forgetting.

\subsection{UMS-in-the-Loop Training}
\label{sec:inloop}

Each modulated projection owns a UMS with column bases
$P,Q\in\mathbb{R}^{d\times R_{\max}}$, a current active rank $r_t$, and a
shared rotation $O\in\mathbb{R}^{d\times d}$.  UMS-in-the-loop is designed to
reduce, rather than eliminate, the mismatch between unconstrained task training
and subsequent subspace encoding.  At task start and at periodic refreshes, the
current spectrum is transformed to the canonical frame,
\begin{equation}
 T_t=O^{\mathsf T}S_tO.
 \label{eq:canonical}
\end{equation}
The implementation removes the component explained by the active left basis,
computes a truncated SVD of the residual, residualizes the resulting right
directions against $Q$, and orthonormalizes both sides by QR.  Appending these
directions to the active columns gives candidate bases
$P_{\mathrm{c}},Q_{\mathrm{c}}\in\mathbb{R}^{d\times k}$, capped by
$R_{\max}$.

Let $A=OP_{\mathrm{c}}$ and $B=OQ_{\mathrm{c}}$.  Restricting the induced
bilinear representation to the current support gives
\begin{equation}
 G_t[m,j]=A[u_m,j]B[v_m,j],
 \qquad G_t\in\mathbb{R}^{n_f\times k}.
 \label{eq:G}
\end{equation}
The coefficients used in the forward pass are
\begin{equation}
 \widetilde{\boldsymbol{\theta}}_t
 =G_t\bigl(G_t^{\mathsf T}\boldsymbol{\theta}_t\bigr),
 \label{eq:inloop_map}
\end{equation}
and are scattered onto $\Omega_t$ before Eq.~\eqref{eq:ifft_modulation}.  The
matrix $G_t$ is cached without gradients and is refreshed at task start and
periodically during training.  Because its support-restricted columns are not
explicitly orthonormal and the code does not apply
$(G_t^{\mathsf T}G_t)^{-1}$, Eq.~\eqref{eq:inloop_map} is a
candidate-subspace-induced linear map, not an orthogonal projector.  Cross
entropy updates the current raw coefficients, bottleneck, and classifier.

\subsection{Post-Task UMS Encoding}
\label{sec:encoding}

The purpose of post-task encoding is to express each learned modulation through
shared bases and a task-specific coordinate that can later be decoded.  At task
completion, the current implementation rebuilds $S_t$ from the
\emph{raw} slice $\boldsymbol{\theta}_t$ and encodes this matrix into UMS.  It
does not fold the final in-loop value
$\widetilde{\boldsymbol{\theta}}_t$ from Eq.~\eqref{eq:inloop_map}; thus the
in-loop mechanism biases learning toward the candidate representation but does
not make training and end-task encoding identical.

Let $r_{t-1}$ be the active rank before encoding task $t$, with $r_0=0$.
For the previously active bases $P_{t-1}=P_{:,1:r_{t-1}}$ and
$Q_{t-1}=Q_{:,1:r_{t-1}}$, the left residual is
\begin{equation}
 U_t=T_t-P_{t-1}(P_{t-1}^{\mathsf T}T_t).
 \label{eq:left_residual}
\end{equation}
A truncated SVD
$U_t\approx U_s\operatorname{diag}(\boldsymbol{\sigma})V_s^{\mathsf T}$
provides paired residual directions.  The right directions are additionally
residualized against $Q_{t-1}$; QR-orthonormalized directions are appended to $P$
and $Q$ when
\begin{equation}
 \mathcal{E}_t=\|U_s\operatorname{diag}(\boldsymbol{\sigma})\|_F^2>\tau
 \label{eq:residual_energy}
\end{equation}
and capacity remains.  In the evaluated configuration,
$R_{\max}=768$, the truncated-SVD budget is 75, and $\tau=10^{-3}$.  With the
updated active rank $r_t$ and bases $P_r=P_{:,1:r_t}$ and
$Q_r=Q_{:,1:r_t}$, task $t$ is encoded as the zero-padded diagonal core
\begin{equation}
 z_t=\operatorname{pad}_{R_{\max}}\!\left[
 \operatorname{diag}\!\left(P_r^{\mathsf T}T_tQ_r\right)\right].
 \label{eq:task_code}
\end{equation}
This is a rank-budgeted diagonal approximation, not an exact factorization of
the complete core.  UMS codes are used for inference, although the present
training/checkpoint implementation also retains the preallocated raw
coefficient tensor.

\subsection{UMS Refinement and Decoding}
\label{sec:multiobjective}

The purpose of refinement is to reduce the discrepancy between the sparse task
spectrum and the spectrum decoded from the diagonal UMS coordinate.  A single
reconstruction term is insufficient for this purpose: it does not identify the
sparse support, expose off-diagonal energy omitted by $z_t$, or control changes
to previously used rotation directions.  We therefore optimize the rotation
and support decoder with a structured objective derived from these encoding
errors.

The shared orientation is parameterized by the Cayley transform
\begin{equation}
 A_{\mathrm{rot}}=W_{\mathrm{rot}}-W_{\mathrm{rot}}^{\mathsf T},\qquad
 O=\left(I-\tfrac{1}{2}A_{\mathrm{rot}}\right)^{-1}
   \left(I+\tfrac{1}{2}A_{\mathrm{rot}}\right),
 \label{eq:cayley}
\end{equation}
which is orthogonal in exact arithmetic whenever the solve is well-defined.
For the current code $z_t$, define the ungated spectral reconstruction and the
support logits as
\begin{align}
 \bar S_t &= OP\operatorname{diag}(z_t)Q^{\mathsf T}O^{\mathsf T},
 \label{eq:ungated_reconstruction}\\
 H_t &= W_L\operatorname{diag}(z_t)W_R^{\mathsf T}+\Gamma.
 \label{eq:support_logits}
\end{align}
Here $W_L,W_R\in\mathbb{R}^{d\times R_{\max}}$ and
$\Gamma\in\mathbb{R}^{d\times d}$ are shared support-decoder parameters.  Let
$\mathcal{I}_t=\{(u,v):S_t[u,v]\neq0\}$,
$Y_t[u,v]=\mathbf{1}[(u,v)\in\mathcal{I}_t]$, and
$C_t=P_r^{\mathsf T}(O^{\mathsf T}S_tO)Q_r$.  For nonempty $\mathcal{I}_t$,
let $m_t$ be the mean of $\bar S_t$ on $\mathcal{I}_t$ and let
$s_t=\sqrt{\operatorname{Var}_{\mathcal{I}_t}(\bar S_t)+10^{-8}}$,
where the variance uses the population denominator.  The code $z_t$ is held
fixed during refinement.  For $t>1$, let $O^{-},P^{-},Q^{-}$ denote the
pre-refinement rotation and previously active basis columns, with
$r_{t-1}$ columns.  The effective refinement objective is
\begin{align}
 \mathcal{L}_{\mathrm{UMS}}
 &=\lambda_{\mathrm{val}}
   \frac{\sum_{(u,v)\in\mathcal{I}_t}(S_t[u,v]-\bar S_t[u,v])^2}
        {|\mathcal{I}_t|}
 +\lambda_{\mathrm{mask}}\operatorname{BCEWithLogits}(H_t,Y_t;w_+)
 \notag\\
 &\quad+\lambda_{\mathrm{mom}}\bigl(m_t^2+(s_t-1)^2\bigr)
 +\lambda_{\mathrm{off}}
   \frac{\|C_t-\operatorname{diag}(\operatorname{diag}C_t)\|_F^2}
        {\max\{r_t(r_t-1),1\}}
 \notag\\
 &\quad+\lambda_{\mathrm{drift}}
   \frac{\|(O-O^{-})P^{-}\|_F^2+\|(O-O^{-})Q^{-}\|_F^2}
        {2r_{t-1}},
 \label{eq:loss_total}
\end{align}
where $w_+=(d^2-n_f)/n_f$ compensates for support imbalance.  The first two
terms penalize value and support-mask errors, respectively; the moment term
regularizes the scale of decoded occupied entries; the off-diagonal term
penalizes canonical-core energy omitted by the diagonal code; and the drift
term penalizes rotation changes along previously active basis directions.  The
value and moment terms are zero when $\mathcal{I}_t$ is empty; the drift term
is zero when there are no previously active columns.  The implemented weights are
$(\lambda_{\mathrm{val}},\lambda_{\mathrm{mask}},\lambda_{\mathrm{mom}},
\lambda_{\mathrm{off}},\lambda_{\mathrm{drift}})
=(1.0,0.1,0.01,0.1,2.0)$.  Refinement uses 15 Adam steps with learning rates
$10^{-4}$ for $W_{\mathrm{rot}}$ and $10^{-3}$ for the support decoder, after
which the current task is re-encoded without further basis expansion.

The implementation also forms a nominal drift penalty for the previously
active columns of $W_L$ and $W_R$.  Its reference tensors are detached views,
however, rather than independent copies: they share storage with the current
parameters and change when Adam updates them.  Consequently, their difference
from the current columns and the resulting penalty and gradient are zero.
Equation~\eqref{eq:loss_total} therefore states the effective drift term.  The
anchor does not constrain $\Gamma$, re-encode all previous tasks, or establish
invariant old-task readouts.  Numerical orthogonality of the Cayley rotation
can be assessed using the saved diagnostics.

At readout, the support decoder produces a hard gate and the task modulation is
\begin{equation}
 M_t=\mathbf{1}[\sigma(H_t)\geq0.5],\qquad
 \widehat{\Delta W}_t
 =\alpha\operatorname{Re}\!\left(
 \mathcal{F}^{-1}_2(\bar S_t\odot M_t)\right).
 \label{eq:ums_readout}
\end{equation}
The gate sets entries predicted as off-support to zero.  This decoding rule is
approximate and does not establish exact reconstruction or absence of
inter-task interactions.

\subsection{Dataset-Level Mahalanobis Routing and Inference}
\label{sec:routing}

The router selects which stored UMS coordinate and task head should be decoded
for a grouped query set.  We apply the Gaussian-memory expert-selection
principle used by LEAR~\cite{yu2025lear}.  At the end of task $t$, the complete
training loader is passed through $f_{\mathrm{route}}$.
If $h_0(x)\in\mathbb{R}^{768}$ is its
class token and $u_t(x)=b_t(h_0(x))\in\mathbb{R}^{500}$, the stored statistics
are
\begin{equation}
 \mu_t=\frac{1}{N_t}\sum_i u_t(x_i^t),\qquad
 \Sigma_t=\operatorname{Cov}\{u_t(x_i^t)\}+10^{-4}I.
 \label{eq:gaussian_memory}
\end{equation}
The diagonal term provides numerical regularization.

The main protocol receives an unlabeled query loader $\mathcal{Q}$ whose
samples are already grouped by source dataset/domain.  To match the evaluator,
write $\mathcal{Q}=\{\mathcal{B}_1,\ldots,\mathcal{B}_J\}$ for its mini-batches.
For candidate task $k$, the evaluator first computes the mean of the square-root
Mahalanobis distances within each batch and then accumulates these batch means
with batch-size weights:
\begin{align}
 d_k(x)&=\sqrt{(u_k(x)-\mu_k)^{\mathsf T}\Sigma_k^{-1}
 (u_k(x)-\mu_k)},\notag\\
 D_{k,j}&=\frac{1}{|\mathcal{B}_j|}\sum_{x\in\mathcal{B}_j}d_k(x),\qquad
 \bar D_k(\mathcal{Q})=
 \frac{\sum_{j=1}^{J}|\mathcal{B}_j|D_{k,j}}
      {\sum_{j=1}^{J}|\mathcal{B}_j|},\notag\\
 k^\star&=\arg\min_k\bar D_k(\mathcal{Q}).
 \label{eq:dataset_router}
\end{align}
Here labels are not used in the distance calculation.  The final $k^\star$
after all batches is assigned to the complete loader; the evaluator does not
route each image independently.  The 24 UMS modules decode $z_{k^\star}$ using
Eq.~\eqref{eq:ums_readout}, and prediction is
\begin{equation}
 \widehat y=c_{k^\star}\!\left(
 b_{k^\star}\bigl(f_{\mathrm{loc}}(x;\widehat{\Delta W}_{k^\star})\bigr)
 \right).
 \label{eq:prediction}
\end{equation}

\subsection{Storage and Computational Characteristics}
\label{sec:storage}

One UMS module stores the fixed shared state
$P,Q,W_L,W_R,W_{\mathrm{rot}},\Gamma$, amounting to
$4dR_{\max}+2d^2$ scalars; with $d=R_{\max}=768$, the 24 modules contain
$84{,}934{,}656$ such scalars.  The UMS code itself grows by
$24R_{\max}=18{,}432$ scalars per task, but the complete task-dependent state
also includes the bottleneck, classifier, and 500-dimensional Gaussian
statistics.  Moreover, the present checkpoint retains the preallocated raw
coefficient tensor.  Hence $n_f/R_{\max}\approx7.68$ characterizes only the
code-to-spectrum ratio for one module, not end-to-end model storage.  Training
and consolidation additionally invoke inverse FFTs, dense SVD/QR operations,
Cayley solves, and support decoding.  Total checkpoint size, consolidation
time, and inference latency are therefore measured empirically in the
experiments rather than inferred from the dimension of $z_t$ alone.
